\setcounter{definition}{0}
\setcounter{theorem}{0}
\subsection{Számelmélet}
\subsubsection{Maradékos osztás}
\begin{theorem}
	Legyen $\mathbb{K} \in \{ \mathbb{C, R, Q, Z}_p \}$ és
	$f, g \in \mathbb{K}[x], g \neq 0$.
	Ekkor egyértelműen léteznek olyan $q, r \in \mathbb{K}[x]$ polinomok, hogy
	\begin{equation*}
		f = g \cdot q + r \quad \deg r < \deg g
	\end{equation*}
\end{theorem}
\begin{proof}
	\begin{gather*}
		f = c_n \cdot x^n + \dots + c_0 \quad
		g = d_m \cdot x^m + \dots + d_0 \quad c_n, d_m \neq 0 \\
		\deg f < \deg g \implies f = g \cdot 0 + f \implies q = 0 \land r = f \\
	\end{gather*}
	\begin{multline*}
		\deg f \geq \deg g \implies n \geq m \implies \\
		f' = f - \frac{c_n}{d_m} \cdot x^{n - m} \cdot g = g \cdot q' + r' \quad
		\deg r' < \deg g
	\end{multline*}
	\begin{equation*}
		f = g \cdot
		\left( q' + \frac{c_n}{d_m} \cdot x^{n - m} \right) + r' \implies
		q = q' + \frac{c_n}{d_m} \cdot x^{n - m} \land r = r'
	\end{equation*}
	\begin{gather*}
		f = g \cdot q + r = g \cdot q' + r' \quad \deg r, \deg r' < \deg g \\
		g \cdot (q - q') = r' - r
	\end{gather*}
	\begin{multline*}
		q \neq q' \implies \deg(q - q') > 0 \implies
		\deg g \cdot (q - q') \geq \deg g \implies \\
		\deg r' - r \geq \deg g
		\tag{Ellentmondás, mert $\deg r, \deg r' < \deg g$}
	\end{multline*}
\end{proof}
\begin{remark*}
	A tétel $\mathbb{Z}[x]$-ben és $\mathbb{Z}_8[x]$-ben igaz.
\end{remark*}
\subsubsection{Legnagyobb közös osztó}
\begin{definition}
	Legyen $f, g \in \mathbb{K}[x]$ polinomok.
	A $h \in \mathbb{K}[x]$ polinom az $f$ és $g$ legnagyobb közös osztója, ha
	\begin{itemize}
		\item $h \mid f$ és $h \mid g$
		\item minden $k \in \mathbb{K}[x]$ esetén,
		ha $k \mid f$ és $k \mid g$, akkor $k \mid h$
		\item $h$ főegyütthatója $1$
	\end{itemize}
\end{definition}
\begin{theorem}
	Bármely két $f, g$ polinomnak létezik legnagyobb közös osztója,
	és az meghatározható az euklideszi algoritmussal.
\end{theorem}
\begin{proof}
	\begin{align*}
		\deg f, \deg g \geq 1 \\
		f = g \cdot q_1 + r_1 & \quad \deg r_1 < \deg g \\
		g = r_1 \cdot q_2 + r_2 & \quad \deg r_2 < \deg r_1 \\
		r_1 = r_2 \cdot q_3 + r_3 & \quad \deg r_3 < \deg r_2 \\
		\dots \\
		r_{\ell - 2} = r_{\ell - 1} \cdot q_\ell + r_\ell & \quad \deg r_\ell < \deg r_{\ell - 1} \\
		r_{\ell - 1} = r_\ell \cdot q_{\ell + 1}
	\end{align*}
	Ekkor $(f, g) = \frac{1}{c} \cdot r_\ell$,
	ahol $c$ az $r_\ell$ polinom főegyütthatója.
\end{proof}
